% META: {"id": "1NSI-TYPES-BASE-RE-C2-CORRIGE", "chapitre": "1NSI-TYPES-BASE", "type_objet": "corrige", "exercice_ref": "1NSI-TYPES-BASE-RE-C2", "capacites": ["P-BASE-02"], "capacites_codes": ["C2"], "mode_creation": "ex_nihilo", "sources_inspiration": [], "status": "needs_review"}
\begin{corrige}{1NSI-TYPES-BASE-RE-C2}
\begin{enumerate}
  \item $2^6 = \mathbf{64}$ valeurs distinctes, de $0$ à $\mathbf{63}$. La plus grande n'est
        pas $64$ : $0$ occupe déjà une des soixante-quatre places.
  \item Pour $100$ : $2^6 = 64 \leq 100 < 128 = 2^7$, il faut donc $\mathbf{7}$ bits.
        Pour $128$ : $2^7 = 128 \leq 128 < 256$, il faut $\mathbf{8}$ bits — l'inégalité est
        stricte à droite, et $128$ ne s'écrit pas sur $7$ bits, qui s'arrêtent à $127$.
  \item Il a calculé $2^7$, c'est-à-dire le \textbf{nombre de valeurs} représentables, et non
        le nombre de bits. Les deux grandeurs se confondent facilement ; l'une est un exposant,
        l'autre une puissance.
\end{enumerate}
\end{corrige}
